\documentclass[11pt]{article}

\usepackage[margin=1in]{geometry}
\usepackage{amsmath, amssymb, mathtools}
\usepackage{xcolor}
\usepackage{enumitem}
\usepackage{array}
\usepackage{hyperref}

\hypersetup{
  colorlinks=true,
  linkcolor=blue,
  urlcolor=blue,
  citecolor=blue
}

% -----------------------------------------------------------------------------
% Basic notation
% -----------------------------------------------------------------------------

\newcommand{\R}{\mathbb{R}}
\newcommand{\N}{\mathbb{N}}
\newcommand{\E}{\mathbb{E}}
\newcommand{\Prob}{\mathbb{P}}
\newcommand{\one}{\mathbf{1}}
\newcommand{\eps}{\varepsilon}

\newcommand{\Dir}{\operatorname{Dirichlet}}
\newcommand{\softmax}{\operatorname{softmax}}
\newcommand{\softplus}{\operatorname{softplus}}
\newcommand{\argmax}{\operatorname*{arg\,max}}
\newcommand{\argmin}{\operatorname*{arg\,min}}
\newcommand{\mean}{\operatorname{mean}}
\newcommand{\stopgrad}{\operatorname{stopgrad}}
\newcommand{\KL}{D_{\mathrm{KL}}}
\newcommand{\DirNLL}{\operatorname{DirNLL}}

\newcommand{\Forest}{\mathcal{F}}
\newcommand{\Tree}{\mathcal{T}}
\newcommand{\Path}{\mathcal{P}}
\newcommand{\Legal}{\mathcal{A}_{\mathrm{legal}}}
\newcommand{\Actions}{\mathcal{A}}
\newcommand{\Z}{\mathcal{Z}}
\newcommand{\Interior}{\mathcal{I}}

\newcommand{\Flip}{\operatorname{flip}}
\newcommand{\Align}{\operatorname{align}}
\newcommand{\Utility}{U}
\newcommand{\Child}{\operatorname{child}}
\newcommand{\Solved}{\operatorname{solved}}
\newcommand{\Cat}{\operatorname{cat}}
\newcommand{\Posterior}{\operatorname{post}}
\newcommand{\EmbedCat}{\operatorname{EmbedCat}}
\newcommand{\SelectDraw}{\operatorname{SelectDraw}}
\newcommand{\SelectDrawDistribution}{\operatorname{SelectDrawDistribution}}

\newcommand{\DecisionEdgeValue}{\operatorname{DecisionEdgeValue}}
\newcommand{\DecisionEdgePosterior}{\operatorname{DecisionEdgePosterior}}
\newcommand{\ChildPosteriorSnapshot}{\operatorname{ChildPosteriorSnapshot}}
\newcommand{\SolveDecision}{\operatorname{SolveDecision}}
\newcommand{\ThompsonSelect}{\operatorname{ThompsonSelect}}
\newcommand{\PosteriorBestPolicyTarget}{\operatorname{PosteriorBestPolicyTarget}}
\newcommand{\RecomputeCache}{\operatorname{RecomputeCache}}
\newcommand{\BackupPath}{\operatorname{BackupPath}}
\newcommand{\NextTransitionRequest}{\operatorname{NextTransitionRequest}}
\newcommand{\RequestTransitions}{\operatorname{RequestTransitions}}
\newcommand{\SubmitTransitions}{\operatorname{SubmitTransitions}}
\newcommand{\FinishSearch}{\operatorname{FinishSearch}}
\newcommand{\MakeTargets}{\operatorname{MakeTargets}}
\newcommand{\AdvanceRoot}{\operatorname{AdvanceRoot}}

% -----------------------------------------------------------------------------
% Lightweight algorithm box
% -----------------------------------------------------------------------------

\newenvironment{algorithmblock}[1]
{
  \par\medskip
  \noindent\hrule\par\smallskip
  \noindent\textbf{#1}
  \par\smallskip
}
{
  \par\smallskip\noindent\hrule\par\medskip
}

\newcommand{\Input}{\textbf{Input: }}
\newcommand{\Output}{\textbf{Output: }}
\newcommand{\Require}{\textbf{Require: }}
\newcommand{\Return}{\textbf{return }}
\newcommand{\If}{\textbf{if }}
\newcommand{\Then}{\textbf{ then }}
\newcommand{\Else}{\textbf{else }}
\newcommand{\EndIf}{\textbf{end if}}
\newcommand{\For}{\textbf{for }}
\newcommand{\EndFor}{\textbf{end for}}
\newcommand{\While}{\textbf{while }}
\newcommand{\EndWhile}{\textbf{end while}}
\newcommand{\Do}{\textbf{ do }}
\newcommand{\Break}{\textbf{break}}
\newcommand{\Continue}{\textbf{continue}}

\title{Dirichlet-Q AlphaZero: Fused Sparse Posterior Search Algorithms}
\author{}
\date{Merged draft -- categorical Dirichlet density NLL}

\begin{document}

\maketitle

\section*{1. Purpose and boundary}

This document specifies the search-side algorithms for Dirichlet-Q AlphaZero in
an implementation where Rust owns the mutable posterior forest and Python owns
environment stepping, legal-action discovery, terminal outcomes, and neural
network evaluation.

\[
\text{Rust owns the mutable posterior forest.}
\]
\[
\text{Python owns transition/evaluation work.}
\]

Rust never calls an environment step function and never contains a game-specific
adapter. Python interacts with Rust through opaque state handles, compact
valid-action tensors, transition request batches, transition submission batches,
and final exported arrays.

The central mathematical object is still a posterior search tree. Each
posterior-valued edge carries a WDL Dirichlet belief. Search refines these
beliefs and trains the policy head toward the posterior probability that each
action is optimal.

The main implementation change from the older algorithm is that tree-internal
parallel repair is removed from the MVP. The MVP invariant is

\[
\#\{\text{pending transition requests in one tree}\}\le 1.
\]

Parallelism comes from batching many trees and roots. A future extension allowing
multiple workers per tree is given in Section 13. That extension uses edge-level
reservation and keeps pending work out of the posterior math.

The MVP supports two structural node kinds:

\[
\texttt{Decision},
\qquad
\texttt{Terminal}.
\]

Decision nodes may additionally become \emph{categorical WDL solved} nodes. In
this document, categorical means an exact solved WDL category in
\(\{L,D,W\}\) carried natively by the posterior tree. It does not mean a chance
node or stochastic environment node. Once a decision node becomes categorical WDL
solved, search stops below that node. Its value is propagated as a native
categorical value, not as a Dirichlet target alpha.

The posterior-only special case is recovered by removing categorical solved tags
and terminal category tags, so every edge value is \(\Posterior(\alpha)\). In
that case the selection, cache, policy target, Q target, and value target reduce
to the old posterior tree formulas.

\section*{2. WDL notation, perspective, and network outputs}

The WDL outcome space is

\[
\Z=\{L,D,W\}.
\]

A WDL probability vector is

\[
\phi=(\phi_L,\phi_D,\phi_W)\in\Delta^2.
\]

The scalar utility used everywhere is

\[
\Utility(\phi)=\phi_W-\phi_L.
\]

For an exact categorical outcome \(z\in\Z\),

\[
\Utility(W)=1,
\qquad
\Utility(D)=0,
\qquad
\Utility(L)=-1.
\]

For a positive Dirichlet vector

\[
\alpha=(\alpha_L,\alpha_D,\alpha_W)\in\R_{>0}^3,
\]

the predictive mean is

\[
\mean(\alpha)=\frac{\alpha}{\alpha_0},
\qquad
\alpha_0=\sum_{z\in\Z}\alpha_z.
\]

For ordinary two-player zero-sum perspective changes,

\[
\Flip(L)=W,
\qquad
\Flip(D)=D,
\qquad
\Flip(W)=L,
\]

and for Dirichlet parameters,

\[
\Flip(\alpha_L,\alpha_D,\alpha_W)=(\alpha_W,\alpha_D,\alpha_L).
\]

All submitted WDL alphas and categorical outcomes are expressed from the
submitted state's current-player perspective. Rust stores the current-player id
for each node and aligns child values to the parent perspective:

\[
\Align_{v\leftarrow u}(x)=
\begin{cases}
x,&\text{if } \texttt{current\_player}(v)=\texttt{current\_player}(u),\\
\Flip(x),&\text{otherwise.}
\end{cases}
\]

For a categorical value \((z,d)\), where \(d\) is distance in plies to the
categorical result, alignment flips only the WDL label:

\[
\Align_{v\leftarrow u}(z,d)=
\left(\Align_{v\leftarrow u}(z),d\right).
\]

At state \(s\), the network snapshot used for search, denoted \(\theta^-\),
predicts

\[
\ell_{\theta^-}(s,a),
\qquad
\alpha_{\theta^-}^V(s),
\qquad
\alpha_{\theta^-}^Q(s,a).
\]

The policy head is

\[
\pi_{\theta^-}(a\mid s)
=
\softmax(\ell_{\theta^-}(s,\cdot))_a.
\]

The value head represents a WDL belief over the state value,

\[
\phi_s^V\sim\Dir(\alpha_{\theta^-}^V(s)),
\]

and the Q head represents a WDL belief over each state-action value,

\[
\phi_{s,a}^Q\sim\Dir(\alpha_{\theta^-}^Q(s,a)).
\]

Both Dirichlet heads use a mean-concentration parameterization. The model emits
WDL mean logits \(r_\theta\) and one scalar concentration logit \(t_\theta\):

\[
\bar\phi_\theta=\softmax(r_\theta),
\qquad
\alpha_{\theta,0}=\softplus(t_\theta)^2,
\qquad
\alpha_\theta=\alpha_{\theta,0}\bar\phi_\theta.
\]

There is no additive base Dirichlet prior in the network parameterization. The
strict positivity of \(\alpha_\theta\) is enforced by the positive concentration
and strictly positive softmax mean.

\section*{3. Sparse posterior forest state}

A forest \(\Forest\) contains one or more active posterior trees. A tree is
searched from its current root. Each decision node stores compact per-action
arrays of length

\[
K_v=|\Legal(s_v)|.
\]

The compact index \(i\in\{0,\ldots,K_v-1\}\) is distinct from the global action
id

\[
a_i=\texttt{legal\_actions}_v[i]
\in\{0,\ldots,A-1\},
\]

where \(A=\texttt{action\_size}\) is only a validation bound for action ids.

For a decision node \(v\), Rust stores:

\[
\begin{array}{ll}
\texttt{state}_v              & \text{opaque Python state handle},\\
\texttt{observation}_v        & \text{observation tensor for export},\\
\texttt{current\_player}_v    & \text{perspective id},\\
\texttt{legal\_actions}_v[i]  & \text{global action id for compact edge }i,\\
\ell_v[i]                    & \text{policy logit for compact edge }i,\\
\alpha_v^V                   & \text{value Dirichlet alpha in }\R_{>0}^3,\\
\alpha_v^Q[i]                & \text{Q Dirichlet alpha for compact edge }i,\\
\Child(v,i)                  & \text{optional child node id},\\
B_{v,i}                      & \text{optional posterior edge snapshot},\\
R_{v,i}                      & \text{completed-result count represented by edge }i,\\
\Cat_{v,i}                   & \text{optional categorical edge value }(z,d),\\
C^V_v                        & \text{optional unsolved state posterior cache},\\
N^\downarrow_v               & \text{completed-result count below }v,\\
\Solved(v)                   & \text{optional solved node value }(z,d,i^\star).
\end{array}
\]

The invariant is

\[
K_v
=|\texttt{legal\_actions}_v|
=|\ell_v|
=|\alpha_v^Q|
=|\Child(v,\cdot)|
=|\Cat_{v,\cdot}|.
\]

No dense policy, Q, child, categorical, or edge array of length \(A\) is stored
on a node.

A terminal node stores an opaque state handle if needed for bookkeeping, the
current-player id at the submitted terminal state, and an intrinsic terminal
category

\[
\Cat_{\mathrm{terminal}}(u)=(z,0).
\]

The intrinsic distance is zero at the terminal node itself. The parent edge that
moves into this terminal node exposes distance one.

\subsection*{3.1 Edge-state fields}

In the MVP, a tree has at most one pending transition request, so each edge can
be represented by simple optional child and optional edge-value fields. For
clarity, an edge is conceptually in one of these states:

\[
\texttt{Empty},
\qquad
\texttt{Pending(token)},
\qquad
\texttt{Expanded(child)},
\qquad
\texttt{Categorical}(z,d).
\]

For the single-pending MVP the \texttt{Pending} state can be stored in the tree's
single pending-request record. For the future multi-worker extension it becomes
an explicit per-edge state; see Section 13.

The fields \(B_{v,i}\), \(R_{v,i}\), and \(C^V_v\) are replacement snapshots.
They are recomputed from current child information and then installed. They are
not additive evidence accumulators.

\section*{4. Python boundary}

\subsection*{4.1 Root insertion}

Python evaluates roots before adding them to Rust. Root insertion receives
compact valid-action rows:

\[
\begin{array}{ll}
\texttt{root\_states}       & [B]\ \text{opaque state handles},\\
\texttt{observations}       & [B,\ldots],\\
\texttt{current\_players}   & [B],\\
\texttt{action\_offsets}    & [B+1],\\
\texttt{legal\_actions}     & [M],\\
\texttt{policy\_logits}     & [M],\\
\texttt{value\_alpha}       & [B,3],\\
\texttt{q\_alpha}           & [M,3].
\end{array}
\]

For row \(r\), valid actions occupy

\[
\texttt{start}=\texttt{action\_offsets}[r],
\qquad
\texttt{end}=\texttt{action\_offsets}[r+1].
\]

The row has \(K_r=\texttt{end}-\texttt{start}\) compact actions. Rust validates
that all action ids are in \([0,A)\), that all alpha entries are positive, and
that the compact arrays have matching lengths.

\subsection*{4.2 Transition requests}

Rust traversal requests fused transition/evaluation work from Python:

\[
\RequestTransitions(\Forest,\texttt{max\_batch\_size},\texttt{pad\_to})
\rightarrow
\texttt{TransitionBatch}.
\]

A transition batch contains:

\[
\begin{array}{ll}
\texttt{token}          & \text{opaque submission token},\\
\texttt{tree\_ids}      & [P],\\
\texttt{parent\_states} & [P]\ \text{opaque state handles},\\
\texttt{actions}        & [P]\ \text{global action ids},\\
\texttt{active\_mask}   & [P]\ \text{true for real rows, false for padding}.
\end{array}
\]

Padding rows repeat valid active rows only to make JAX shapes static. Padding
rows must be ignored by Rust after shape validation.

Python performs exactly one fused step for the active rows,

\[
\texttt{child\_states}=\texttt{env.step(parent\_states,actions)},
\]

and evaluates the resulting nonterminal child observations with the neural
network. Python then compacts each nonterminal child's dense model outputs to
valid actions before submission.

\subsection*{4.3 Transition submission}

The submission API is

\[
\SubmitTransitions(\Forest,\texttt{token},\texttt{batch outputs}).
\]

For each active row, Python supplies:

\[
\begin{array}{ll}
\texttt{child\_states}       & [P]\ \text{opaque state handles},\\
\texttt{observations}        & [P,\ldots]\ \text{for nonterminal rows},\\
\texttt{current\_players}    & [P],\\
\texttt{terminated}          & [P]\ \text{boolean},\\
\texttt{terminal\_outcomes}  & [P]\ \text{WDL label for terminal rows},\\
\texttt{action\_offsets}     & [P+1]\ \text{compact nonterminal segments},\\
\texttt{legal\_actions}      & [M],\\
\texttt{policy\_logits}      & [M],\\
\texttt{value\_alpha}        & [P,3]\ \text{for nonterminal rows},\\
\texttt{q\_alpha}            & [M,3].
\end{array}
\]

For a terminal row, \(\texttt{terminal\_outcomes}[r]\in\Z\) is authoritative.
The terminal child stores intrinsic value \((z,0)\). The parent edge to that
terminal child exposes \((\Align_{v\leftarrow u}(z),1)\). Terminal outcomes are
not converted into Dirichlet target alphas.

\section*{5. Edge values and posterior fallback}

An edge value is a tagged value:

\[
E_{v,i}\in
\left\{
\Posterior(\alpha),\ \Cat(z,d)
\right\}.
\]

\(\Posterior(\alpha)\) means the edge is represented by a WDL Dirichlet alpha.
\(\Cat(z,d)\) means the edge is exactly solved as WDL category \(z\) at distance
\(d\) from the parent node. A categorical edge is not represented as a Dirichlet
vector inside the posterior tree.

\subsection*{5.1 Posterior fallback}

For a decision node \(v\) and compact edge \(i\), the posterior fallback is

\[
\DecisionEdgePosterior(\Tree,v,i)=
\begin{cases}
B_{v,i},&B_{v,i}\ne\bot,\\[2pt]
\Align_{v\leftarrow u}(C^V_u),
&\begin{array}{l}
u=\Child(v,i),\\
u\text{ has an unsolved posterior cache},
\end{array}\\[8pt]
\Align_{v\leftarrow u}(\alpha_u^V),
&\begin{array}{l}
u=\Child(v,i),\\
u\text{ is an expanded unsolved decision node},
\end{array}\\[8pt]
\alpha_v^Q[i],&\text{otherwise}.
\end{cases}
\]

This function returns only posterior Dirichlet alphas. It is used when an edge
has no categorical value.

\subsection*{5.2 Full edge value}

The full edge value checks categorical information before posterior fallback:

\[
\alpha^{\mathrm{edge}}_{v,i}
=
\DecisionEdgePosterior(\Tree,v,i),
\]
\[
\DecisionEdgeValue(\Tree,v,i)=
\begin{cases}
\Cat(z,d),&\Cat_{v,i}=(z,d),\\[2pt]
\Cat(\Align_{v\leftarrow u}(z),1),
&\begin{array}{l}
u=\Child(v,i),\\
u\text{ terminal with }\Cat_{\mathrm{terminal}}(u)=(z,0),
\end{array}\\[8pt]
\Cat(\Align_{v\leftarrow u}(z),d+1),
&\begin{array}{l}
u=\Child(v,i),\\
\Solved(u)=(z,d,\cdot),
\end{array}\\[8pt]
\Posterior(\alpha^{\mathrm{edge}}_{v,i}),&\text{otherwise}.
\end{cases}
\]


The \(+1\) distance accounts for taking action \(i\) from \(v\) to its child.
This convention is important: a terminal node has intrinsic distance zero, but
an edge that reaches that terminal in one move has edge distance one.

\subsection*{5.3 Explicit child posterior snapshot}

Backup must not define a new posterior snapshot in terms of
\(\DecisionEdgePosterior(\Tree,v,i)\), because that function first checks the old
snapshot \(B_{v,i}\). Instead, when the child \(u=\Child(v,i)\) is unsolved and
posterior-valued, the replacement snapshot is

\[
\ChildPosteriorSnapshot(\Tree,v,i,u)=
\begin{cases}
\Align_{v\leftarrow u}(C^V_u),&u\text{ has a recomputed posterior cache},\\[2pt]
\Align_{v\leftarrow u}(\alpha_u^V),
&u\text{ is an expanded unsolved decision node},\\[2pt]
\alpha_v^Q[i],&\text{otherwise}.
\end{cases}
\]

Backup installs

\[
B_{v,i}\gets\ChildPosteriorSnapshot(\Tree,v,i,u),
\]

not a self-referential call to \(\DecisionEdgePosterior\).

\section*{6. Native categorical WDL solving}

Categorical WDL solving is separate from posterior backup. It is exact
minimax-style propagation over known solved children.

For a decision node \(v\), define

\[
S_v=\{(i,z_i,d_i):\DecisionEdgeValue(\Tree,v,i)=\Cat(z_i,d_i)\}.
\]

The solver is:

\begin{algorithmblock}{Algorithm 1: SolveDecision\((\Tree,v,\texttt{draw\_mode})\)}

\Input decision node \(v\) with \(K_v\) compact legal actions

\Output \(\bot\), or solved value \((z,d,i^\star)\)

\begin{enumerate}[leftmargin=2em]
  \item Recompute categorical edge values from categorical children.

  \item \If there exists \((i,W,d)\in S_v\) \Then
  \[
  i^\star\gets\argmin_{(i,W,d)\in S_v} d.
  \]
  \[
  \Return (W,d_{i^\star},i^\star).
  \]
  \EndIf

  \item \If \(|S_v|<K_v\) \Then \Return \(\bot\). \EndIf

  \item \If every element of \(S_v\) has label \(L\) \Then
  \[
  i^\star\gets\argmax_{(i,L,d)\in S_v} d.
  \]
  \[
  \Return (L,d_{i^\star},i^\star).
  \]
  \EndIf

  \item \If every element of \(S_v\) has label \(D\) or \(L\), and at least one
  element has label \(D\), \Then
  \[
  i^\star\gets\SelectDraw(\{(i,D,d)\in S_v\},\texttt{draw\_mode}).
  \]
  \[
  \Return (D,d_{i^\star},i^\star).
  \]
  \EndIf

  \item \Return \(\bot\).
\end{enumerate}

\end{algorithmblock}

Consequences:

\begin{enumerate}[leftmargin=2em]
  \item A single known winning child solves the node immediately as a win.
  \item The chosen winning action is the shortest win.
  \item A loss requires every legal action to be a known loss.
  \item The chosen losing action is the longest loss.
  \item A draw requires every legal action to be a known draw or known loss,
  with at least one draw.
  \item The chosen draw action is controlled by an engine configuration parameter
  \(\texttt{draw\_mode}\).
  \item If any legal child is still unknown and no win is known, the node
  remains searchable.
\end{enumerate}

\(\SelectDraw\) is an implementation-defined enum selected at engine
construction, not a Python callback during search. Examples include
\(\texttt{shortest\_draw}\), \(\texttt{longest\_draw}\),
\(\texttt{policy\_prior}\), and \(\texttt{uniform}\).

\section*{7. Thompson selection and posterior-best policy targets}

For an unsolved decision node, Thompson selection samples or reads a score for
every compact edge:

\[
X_i=
\begin{cases}
\Utility(z_i),&\DecisionEdgeValue(\Tree,v,i)=\Cat(z_i,d_i),\\[2pt]
\Utility(\phi_i),&\DecisionEdgeValue(\Tree,v,i)=\Posterior(\alpha_i),\
\quad \phi_i\sim\Dir(\alpha_i).
\end{cases}
\]

The selected compact edge is

\[
i^\star=\argmax_i X_i.
\]

Ties between categorical edges are broken by the same distance rules used by the
solver: shortest win, longest loss, and \(\texttt{draw\_mode}\) for draws. The
returned action is the global action id

\[
a^\star=\texttt{legal\_actions}_v[i^\star].
\]

The posterior-best policy target for an unsolved node is the Monte Carlo
probability that each compact edge is selected under this mixed
posterior/categorical scoring rule:

\[
\pi_{\mathrm{search}}^v(i)
=
\Prob\left(
  i=\argmax_j X_j
\right).
\]

It is estimated with \(M_\pi\) independent samples:

\[
\hat\pi_{\mathrm{search}}^v(i)
=
\frac{1}{M_\pi}
\sum_{m=1}^{M_\pi}
\one\left[i=\argmax_j X_j^{(m)}\right].
\]

For a solved categorical node, no Monte Carlo sampling is needed. The policy
target is the configured solved-action distribution:

\[
\pi_{\mathrm{search}}^v(i)=
\begin{cases}
1,&i=i^\star\ \text{for win/loss solved nodes},\\
\SelectDrawDistribution(i),&\Solved(v)=(D,d,i^\star),\\
0,&\text{otherwise}.
\end{cases}
\]

Unlike standard AlphaZero, this policy target is not a visit-count distribution.
For unsolved nodes it is the posterior probability that each move is optimal.

\section*{8. Search-weighted posterior cache for unsolved nodes}

Only unsolved posterior-valued decision nodes maintain a Dirichlet state cache
\(C^V_v\). Categorical solved nodes expose \(\Solved(v)\) instead. The state
cache is a full replacement recomputation from current edge values. It is not an
additive evidence update.

Categorical edge values remain native tags for solving, selection, finish, and
export. However, an unsolved node may have a mixture of posterior and categorical
outgoing edges. To compute a Dirichlet state cache for such an unsolved node, use
a cache-only embedding of exact WDL labels:

\[
\EmbedCat(z)
=
\eps_{\mathrm{cat-cache}}\one
+
\kappa_{\mathrm{cat-cache}}e_z,
\qquad
\eps_{\mathrm{cat-cache}}>0,
\quad
\kappa_{\mathrm{cat-cache}}>0.
\]

This embedding is never exported as a categorical Q or V target and never
replaces the native categorical tag. It exists only so that mixed
categorical/posterior outgoing edges can produce an unsolved node's Dirichlet
cache.

For each compact edge define the cache vector

\[
\widetilde\alpha_{v,i}=
\begin{cases}
\alpha_i,&\DecisionEdgeValue(\Tree,v,i)=\Posterior(\alpha_i),\\[2pt]
\EmbedCat(z_i),&\DecisionEdgeValue(\Tree,v,i)=\Cat(z_i,d_i).
\end{cases}
\]

Let \(R_{v,i}\) be the completed-result count represented by edge \((v,i)\). An
unexpanded fallback edge has \(R_{v,i}=0\). An expanded posterior child has
\(R_{v,i}=1+N^\downarrow_u\), counting the child's own neural evaluation plus
completed work below it. A terminal or solved categorical edge has a positive
count chosen by the implementation, with the default

\[
R_{v,i}=1+N^\downarrow_u
\]

when the categorical value came from a child \(u\), and \(R_{v,i}=1\) for a direct
terminal edge. These counts affect only how much the state cache trusts
completed downstream work; the Dirichlet vectors themselves carry the WDL
beliefs and concentrations.

The local posterior-best policy is computed over all compact edges by the mixed
rule in Section 7:

\[
\hat\pi_{\mathrm{search}}^v
=
\PosteriorBestPolicyTarget
\left(
\{\DecisionEdgeValue(\Tree,v,i)\}_{i=0}^{K_v-1},M_\pi
\right).
\]

Then define the search-weighted edge mixture

\[
E^V_v
=
\sum_{i=0}^{K_v-1}
\hat\pi_{\mathrm{search}}^v(i)\widetilde\alpha_{v,i}.
\]

All legal actions participate. Unexpanded fallback actions are not masked out. If
an unexpanded action has tiny posterior-best probability, it contributes almost
nothing; if it still has large posterior-best probability, the cache should
retain that uncertainty.

The downstream completed-result count is

\[
N^\downarrow_v
=
\sum_{i=0}^{K_v-1} R_{v,i}.
\]

Given a half-trust hyperparameter \(\kappa_N>0\), set

\[
\gamma_v
=
\frac{N^\downarrow_v}{\kappa_N+N^\downarrow_v}.
\]

The cached state posterior is

\[
C^V_v
=
(1-\gamma_v)\alpha_v^V+
\gamma_v E^V_v.
\]

Thus little downstream work leaves \(C^V_v\) close to the neural value prior
\(\alpha_v^V\), while many completed results move \(C^V_v\) toward the
search-weighted posterior over outgoing edges.

\begin{algorithmblock}{Algorithm 2: RecomputeCache\((\Tree,v,M_\pi,\kappa_N)\)}

\Input unsolved expanded decision node \(v\)

\Output \((C^V_v,N^\downarrow_v)\)

\begin{enumerate}[leftmargin=2em]
  \item For every compact edge \(i\), compute
  \[
  E_{v,i}\gets\DecisionEdgeValue(\Tree,v,i).
  \]

  \item Compute the local posterior-best policy
  \[
  \hat\pi_{\mathrm{search}}^v
  \gets
  \PosteriorBestPolicyTarget(\{E_{v,i}\}_{i=0}^{K_v-1},M_\pi).
  \]

  \item Convert each edge to a cache vector
  \[
  \widetilde\alpha_{v,i}\gets
  \begin{cases}
  \alpha_i,&E_{v,i}=\Posterior(\alpha_i),\\
  \EmbedCat(z_i),&E_{v,i}=\Cat(z_i,d_i).
  \end{cases}
  \]

  \item Compute
  \[
  E^V_v\gets
  \sum_{i=0}^{K_v-1}
  \hat\pi_{\mathrm{search}}^v(i)\widetilde\alpha_{v,i}.
  \]

  \item Compute
  \[
  N^\downarrow_v\gets\sum_{i=0}^{K_v-1}R_{v,i},
  \qquad
  \gamma_v\gets\frac{N^\downarrow_v}{\kappa_N+N^\downarrow_v}.
  \]

  \item Install the replacement cache
  \[
  C^V_v\gets(1-\gamma_v)\alpha_v^V+\gamma_v E^V_v.
  \]

  \item \Return \((C^V_v,N^\downarrow_v)\).
\end{enumerate}

\end{algorithmblock}

\section*{9. Traversal and transition scheduling}

Each tree has at most one pending transition request in the MVP. Therefore
duplicate leaf requests and in-tree write races are structurally absent.

\begin{algorithmblock}{Algorithm 3: NextTransitionRequest\((\Tree)\)}

\Output no request, completed simulation, tree done, or one transition request

\begin{enumerate}[leftmargin=2em]
  \item \If the tree already has a pending request \Then \Return no request.
  \EndIf

  \item Start at the current root \(r\), with an empty path \(\Path\).

  \item \While true \Do
  \begin{enumerate}[leftmargin=2em]
    \item \If \(v\) is terminal or \(\Solved(v)\ne\bot\) \Then
    \[
    \BackupPath(\Tree,\Path,v)
    \]
    and \Return completed simulation. \EndIf

    \item Recompute
    \[
    \Solved(v)\gets\SolveDecision(\Tree,v,\texttt{draw\_mode}).
    \]
    \If \(\Solved(v)\ne\bot\) \Then
    \[
    \BackupPath(\Tree,\Path,v)
    \]
    and \Return completed simulation. \EndIf

    \item Select compact edge
    \[
    i\gets\ThompsonSelect(\Tree,v).
    \]

    \item \If \(\Child(v,i)\) exists \Then append \((v,i)\) to \(\Path\), set
    \(v\gets\Child(v,i)\), and continue. \EndIf

    \item Create a pending transition request for
    \[
    (\texttt{state}_v,\texttt{legal\_actions}_v[i]).
    \]
    Store the token, tree generation, parent node id, compact edge index, and
    path. \Return transition request.
  \end{enumerate}
  \EndWhile
\end{enumerate}

\end{algorithmblock}

\begin{algorithmblock}{Algorithm 4: RequestTransitions\((\Forest,\texttt{max\_batch\_size},\texttt{pad\_to})\)}

\begin{enumerate}[leftmargin=2em]
  \item Iterate active trees in round-robin order.
  \item For each tree, call \(\NextTransitionRequest(\Tree)\).
  \item Append at most one active transition row per tree.
  \item Stop when \(\texttt{max\_batch\_size}\) active rows have been collected
  or all active trees are pending or done.
  \item If \(\texttt{pad\_to}\) is set, repeat active rows to produce a static
  padded batch and mark repeated rows inactive.
  \item Return the token, parent state handles, global action ids, tree ids, and
  active mask.
\end{enumerate}

\end{algorithmblock}

\section*{10. Submission, backup, and cache recomputation}

\begin{algorithmblock}{Algorithm 5: SubmitTransitions\((\Forest,\texttt{token},\texttt{outputs})\)}

\begin{enumerate}[leftmargin=2em]
  \item Validate the token and all batch shapes.
  \item Ignore inactive padded rows after validation.
  \item For each active row:
  \begin{enumerate}[leftmargin=2em]
    \item Look up the pending request. If the tree generation has changed,
    discard the row as stale.
    \item Let \((v,i)\) be the pending parent node and compact edge.
    \item If \(\texttt{terminated}\) is true, create a terminal child \(u\) with
    intrinsic categorical value \((z,0)\), where
    \(z=\texttt{terminal\_outcomes}[r]\).
    \item If \(\texttt{terminated}\) is false, create a decision child \(u\) from
    the submitted compact valid-action row, policy logits, value alpha, Q alpha,
    current player, observation, and child state.
    \item Attach the child to \(\Child(v,i)\).
    \item Clear the pending request for the tree.
    \item Run
    \[
    \BackupPath(\Tree,\Path\mathbin{+}[(v,i)],u).
    \]
  \end{enumerate}
\end{enumerate}

\end{algorithmblock}

Backup updates two independent pieces of state:

\begin{enumerate}[leftmargin=2em]
  \item native categorical solve state, via \(\SolveDecision\), and
  \item posterior Dirichlet edge snapshots and posterior caches for unsolved
  nodes.
\end{enumerate}

For a parent edge \((v,i)\) whose child \(u\) is terminal with intrinsic terminal
value \((z,0)\), backup stores the categorical edge value

\[
\Cat_{v,i}\gets(\Align_{v\leftarrow u}(z),1),
\qquad
R_{v,i}\gets 1.
\]

For a parent edge \((v,i)\) whose child \(u\) is solved categorical with
\(\Solved(u)=(z,d,\cdot)\), backup stores

\[
\Cat_{v,i}\gets(\Align_{v\leftarrow u}(z),d+1),
\qquad
R_{v,i}\gets 1+N^\downarrow_u.
\]

No Dirichlet alpha is constructed for either categorical edge value.

For a parent edge \((v,i)\) whose child \(u\) is unsolved posterior-valued,
backup stores a full replacement posterior snapshot

\[
B_{v,i}\gets\ChildPosteriorSnapshot(\Tree,v,i,u),
\qquad
R_{v,i}\gets 1+N^\downarrow_u.
\]

If \(u\) has just been expanded and has no child evidence, then
\(N^\downarrow_u=0\) and \(B_{v,i}=\Align_{v\leftarrow u}(\alpha_u^V)\). The
count one represents the child's own neural evaluation.

\begin{algorithmblock}{Algorithm 6: BackupPath\((\Tree,\Path,u)\)}

\Input completed child or reached node \(u\), with \(\Path=[(v_0,i_0),\ldots,(v_m,i_m)]\)

\Output updated edge snapshots, solved states, and caches along the path

\begin{enumerate}[leftmargin=2em]
  \item Set \(x\gets u\).

  \item Traverse \(\Path\) in reverse order. For each parent edge \((v,i)\) whose
  child is the current \(x\):
  \begin{enumerate}[leftmargin=2em]
    \item If \(x\) is terminal, install
    \[
    \Cat_{v,i}\gets(\Align_{v\leftarrow x}(z_x),1),
    \qquad
    R_{v,i}\gets 1.
    \]

    \item Else if \(\Solved(x)=(z,d,\cdot)\), install
    \[
    \Cat_{v,i}\gets(\Align_{v\leftarrow x}(z),d+1),
    \qquad
    R_{v,i}\gets1+N^\downarrow_x.
    \]

    \item Else \(x\) is an unsolved decision node. Recompute its cache if needed:
    \[
    (C^V_x,N^\downarrow_x)\gets\RecomputeCache(\Tree,x,M_\pi,\kappa_N).
    \]
    Then install
    \[
    B_{v,i}\gets\Align_{v\leftarrow x}(C^V_x),
    \qquad
    R_{v,i}\gets1+N^\downarrow_x.
    \]

    \item Recompute
    \[
    \Solved(v)\gets\SolveDecision(\Tree,v,\texttt{draw\_mode}).
    \]

    \item If \(\Solved(v)=\bot\), recompute the unsolved posterior cache
    \[
    (C^V_v,N^\downarrow_v)\gets\RecomputeCache(\Tree,v,M_\pi,\kappa_N).
    \]

    \item Set \(x\gets v\) and continue upward.
  \end{enumerate}

  \item Return after the root has been updated.
\end{enumerate}

\end{algorithmblock}

This backup rule is intentionally serial and local in the MVP. There are no
dirty nodes, repair tokens, child-cache absorbed versions, or compare-and-swap
node expansion states in the main algorithm. A future multi-worker version may
use local locks and edge reservations, but completed backups must still be
mathematically equivalent to some serial order of these replacement operations.

\section*{11. Finish and root advancement}

\subsection*{11.1 Finish}

\[
\FinishSearch(\Forest,\texttt{tree\_ids},\texttt{commit\_mode})
\]

requires no pending request for the requested trees. Before export, Rust
recomputes caches for all exported roots and interior nodes deterministically
from current edge values.

If the root is solved categorical, finish returns:

\[
\begin{array}{ll}
\texttt{root\_kind}      & \texttt{categorical},\\
\texttt{root\_outcome}   & z,\\
\texttt{root\_distance}  & d,\\
\texttt{action}          & \texttt{legal\_actions}[i^\star],\\
\texttt{policy\_target}  & \text{solved-action distribution}.
\end{array}
\]

If the root is unsolved posterior-valued, finish returns compact sparse
posterior-search arrays:

\[
\begin{array}{ll}
\texttt{action\_offsets} & [B+1],\\
\texttt{legal\_actions}  & [M],\\
\texttt{pi\_search}      & [M],\\
\texttt{root\_q\_kind}   & [M]\ \texttt{posterior or categorical},\\
\texttt{root\_q\_alpha}  & [M,3]\ \text{for posterior entries},\\
\texttt{root\_q\_cat}    & [M]\ \text{for categorical entries},\\
\texttt{actions}         & [B]\ \text{committed global action ids}.
\end{array}
\]

For an unsolved root, \(\texttt{pi\_search}\) is the posterior-best policy target
from Section 7. Commit modes may include:

\[
\texttt{posterior\_sample},
\qquad
\texttt{posterior\_argmax},
\qquad
\texttt{mean\_utility\_argmax}.
\]

The policy training target is unchanged by the commit mode.

\subsection*{11.2 Root advancement}

Root advancement promotes an existing child subtree:

\[
\AdvanceRoot(\Tree,a).
\]

\begin{enumerate}[leftmargin=2em]
  \item Require that the current root is a decision node.
  \item Convert global action id \(a\) to compact edge index \(i\).
  \item Require \(\Child(r,i)\) to exist.
  \item Increment the tree generation so outstanding transition tokens become
  stale.
  \item Promote \(\Child(r,i)\) to the new root and prune unreachable nodes.
  \item Preserve the promoted subtree's posterior data and categorical solved
  state.
\end{enumerate}

Rust must not synthesize a missing child during advancement. If the action was
never expanded, Python must add a fresh root instead.

\section*{12. Training target export and losses}

Training export is sparse and compact over valid actions. It excludes terminal
nodes. It may include both unsolved posterior decision nodes and solved
categorical decision nodes.

For each exported decision node \(v\), append one compact segment:

\[
\begin{array}{ll}
\texttt{observations}        & [N,\ldots],\\
\texttt{action\_offsets}     & [N+1],\\
\texttt{legal\_actions}      & [M],\\
\texttt{policy\_target}      & [M],\\
\texttt{v\_target\_kind}     & [N]\ \texttt{posterior or categorical},\\
\texttt{v\_target\_alpha}    & [N,3]\ \text{for posterior rows},\\
\texttt{v\_target\_cat}      & [N]\ \text{for categorical rows},\\
\texttt{q\_target\_kind}     & [M]\ \texttt{posterior or categorical},\\
\texttt{q\_target\_alpha}    & [M,3]\ \text{for posterior edges},\\
\texttt{q\_target\_cat}      & [M]\ \text{for categorical edges},\\
\texttt{q\_loss\_weight}     & [M],\\
\texttt{row\_mask}           & [N].
\end{array}
\]

For a solved categorical node, the value target is its native WDL category. For
a categorical edge, the Q target is its native WDL category. Neither is exported
as a target Dirichlet alpha. At learner time only, each native category \(z\) is
converted to the smoothed simplex point \(\phi_{\eps_{\mathrm{cat}}}(z)\) for
the Dirichlet density NLL.

For an unsolved posterior-valued node, the value target is the node cache:

\[
\alpha^{V,\mathrm{target}}_v=C^V_v.
\]

For a posterior-valued edge, the Q target is the posterior currently used by
search:

\[
\alpha^{Q,\mathrm{target}}_{v,i}
=
\DecisionEdgePosterior(\Tree,v,i).
\]

The default Q loss weight is the posterior-best policy target:

\[
w_Q(v,i)=\stopgrad(\hat\pi_{\mathrm{search}}^v(i)).
\]

\subsection*{12.1 Policy loss}

The policy loss is cross entropy against the exported posterior-best or solved
action distribution:

\[
\mathcal{L}_\pi(v)
=
-
\sum_{i=0}^{K_v-1}
\stopgrad(\hat\pi_{\mathrm{search}}^v(i))
\log \pi_\theta(\texttt{legal\_actions}_v[i]\mid s_v).
\]

\subsection*{12.2 Posterior Dirichlet losses}

For posterior value and posterior Q targets, train by Dirichlet KL:

\[
\mathcal{L}_V^{\mathrm{post}}(v)
=
\KL\left(
\Dir(\stopgrad(\alpha^{V,\mathrm{target}}_v))
\,\|\,
\Dir(\alpha_\theta^V(s_v))
\right),
\]

and, with \(a_i=\texttt{legal\_actions}_v[i]\),

\[
\mathcal{L}_Q^{\mathrm{post}}(v)
=
\frac{
\sum_{i=0}^{K_v-1}w_Q(v,i)
\KL\left(
\Dir(\stopgrad(\alpha^{Q,\mathrm{target}}_{v,i}))
\,\|\,
\Dir(\alpha_\theta^Q(s_v,a_i))
\right)
}{
\sum_{i=0}^{K_v-1}w_Q(v,i)+\eps
}.
\]

The Dirichlet KL is

\[
\begin{aligned}
\KL(\Dir(\beta)\,\|\,\Dir(\alpha))
&=
\log\Gamma(\beta_0)-\sum_z\log\Gamma(\beta_z)
-
\log\Gamma(\alpha_0)+\sum_z\log\Gamma(\alpha_z)\\
&\quad+
\sum_z(\beta_z-\alpha_z)
\left[\psi(\beta_z)-\psi(\beta_0)\right],
\end{aligned}
\]

where \(\beta_0=\sum_z\beta_z\), \(\alpha_0=\sum_z\alpha_z\), and \(\psi\) is
the digamma function.

\subsection*{12.3 Dirichlet density NLL for categorical neural targets}

The model still outputs positive Dirichlet parameters \(\alpha\) for value and Q
heads. For a categorical WDL target \(z\in\Z\), train the full Dirichlet head by
evaluating the Dirichlet density at a smoothed categorical WDL point.

Define an epsilon-smoothed categorical point
\(\phi_{\eps_{\mathrm{cat}}}(z)\in\Interior\), where
\(\Interior=\{\phi\in\Delta^2: \phi_i>0\ \forall i\}\):

\[
\phi_{\eps_{\mathrm{cat}}}(z)_i=
\begin{cases}
1-\eps_{\mathrm{cat}},&i=z,\\[2pt]
\eps_{\mathrm{cat}}/(|\Z|-1),&i\ne z.
\end{cases}
\]

Since \(|\Z|=3\), the off-target mass is \(\eps_{\mathrm{cat}}/2\). Require

\[
0<\eps_{\mathrm{cat}}<1,
\]

so every logarithm below is finite.

Let \(\alpha_\theta\) be the current prediction of either the value Dirichlet
head or the action Dirichlet-Q head, and define

\[
\alpha_{\theta,0}=\sum_{i\in\Z}\alpha_{\theta,i}.
\]

The categorical neural loss is the Dirichlet density negative log-likelihood at
that smoothed point:

\[
\DirNLL(\alpha_\theta,z;\eps_{\mathrm{cat}})
=
-
\log p_{\Dir(\alpha_\theta)}
\left(
\phi_{\eps_{\mathrm{cat}}}(z)
\right).
\]

Expanding the Dirichlet density gives

\[
\DirNLL(\alpha_\theta,z;\eps_{\mathrm{cat}})
=
-
\log\Gamma(\alpha_{\theta,0})
+
\sum_{i\in\Z}\log\Gamma(\alpha_{\theta,i})
-
\sum_{i\in\Z}
(\alpha_{\theta,i}-1)
\log \phi_{\eps_{\mathrm{cat}}}(z)_i.
\]

The CPU tree still stores categorical targets as native WDL tags. It does not
store \(\phi_{\eps_{\mathrm{cat}}}(z)\) and does not convert categorical targets
into target Dirichlet alphas. The smoothed point is created only inside the
JAX/Python loss code.

This objective is not categorical predictive cross entropy. It is a density NLL
for the Dirichlet distribution over WDL probability vectors, so it trains both
the predictive mean and the concentration. Because it is a density loss, its
numerical value can be negative without indicating an error.

\subsection*{12.4 Combined row loss}

For an exported row \(v\), choose the value loss from the value target kind:

\[
\mathcal{L}_V(v)=
\begin{cases}
\mathcal{L}_V^{\mathrm{post}}(v),
&\texttt{v\_target\_kind}=\texttt{posterior},\\[2pt]
\DirNLL(\alpha_\theta^V(s_v),z_v;\eps_{\mathrm{cat}}),
&\texttt{v\_target\_kind}=\texttt{categorical}.
\end{cases}
\]

For Q, mix posterior and categorical losses edge by edge:

\[
\mathcal{L}_Q(v)=
\frac{
\sum_{i=0}^{K_v-1}w_Q(v,i)\ell_Q(v,i)
}{
\sum_{i=0}^{K_v-1}w_Q(v,i)+\eps
},
\]

where \(k^Q_i=\texttt{q\_target\_kind}[i]\) and

\[
\ell_Q(v,i)=
\begin{cases}
\KL\left(
\Dir(\stopgrad(\alpha^{Q,\mathrm{target}}_{v,i}))
\,\|\,
\Dir(\alpha_\theta^Q(s_v,a_i))
\right),
&k^Q_i=\texttt{posterior},\\[4pt]
\DirNLL(\alpha_\theta^Q(s_v,a_i),z_{v,i};\eps_{\mathrm{cat}}),
&k^Q_i=\texttt{categorical}.
\end{cases}
\]

The total row loss is

\[
\mathcal{L}(v)
=
\mathcal{L}_\pi(v)
+\lambda_V\mathcal{L}_V(v)
+\lambda_Q\mathcal{L}_Q(v),
\]

with row masks applied outside this expression.

\begin{algorithmblock}{Algorithm 7: MakeTargets\((\Forest,\texttt{tree\_ids},M_\pi,\kappa_N)\)}

\begin{enumerate}[leftmargin=2em]
  \item Require no pending transition request for the requested trees.
  \item Recompute caches for exported unsolved nodes from current edge values.
  \item Initialize an empty row list \(\mathcal{D}\gets()\).
  \item For each exported decision node \(v\):
  \begin{enumerate}[leftmargin=2em]
    \item Export its observation and compact legal-action segment.
    \item Export the policy target: posterior-best for unsolved nodes, or
    solved-action distribution for solved categorical nodes.
    \item Export value kind and value target: \(C^V_v\) for unsolved posterior
    nodes, native category for solved categorical nodes.
    \item For each compact edge, export Q kind and Q target: posterior alpha for
    posterior edges, native category for categorical edges.
    \item Export Q loss weight, defaulting to the policy target weight.
  \end{enumerate}
  \item Concatenate rows from all trees, shuffle, pad if needed, and return
  learner arrays.
\end{enumerate}

\end{algorithmblock}

\section*{13. Future extension: multiple workers per tree}

The MVP uses one pending transition request per tree. To allow multiple workers
per tree without returning to the old dirty/repair system, replace the tree-level
pending invariant with an edge-level reservation invariant:

\[
\#\{\text{pending requests for edge }(v,i)\}\le 1,
\qquad
\#\{\text{pending requests in tree }\Tree\}\le W_{\mathrm{tree}}.
\]

Each edge has an explicit concurrent state:

\[
\texttt{Empty}
\rightarrow
\texttt{Pending(token)}
\rightarrow
\texttt{Expanded(child)},
\]

or

\[
\texttt{Empty}
\rightarrow
\texttt{Pending(token)}
\rightarrow
\texttt{Categorical}(z,d).
\]

When a worker selects an empty edge, it attempts an atomic reservation:

\[
\texttt{Empty}\xrightarrow{\mathrm{CAS}}\texttt{Pending(token)}.
\]

If it wins, it emits a transition request. If it loses because another worker
already reserved the edge, it does not emit a duplicate request.

A worker that reaches a pending edge should use a scheduling rule, not a
posterior update. The simplest rule is to temporarily skip pending edges during
selection unless all edges are pending. An alternative is a virtual-loss-like
pending penalty, but it must be stored separately as scheduling-only state, for
example

\[
P^{\mathrm{inflight}}_{v,i},
\]

and must never be added to

\[
B_{v,i},
\qquad
C^V_v,
\qquad
\alpha_v^Q[i],
\qquad
\Cat_{v,i}.
\]

Pending requests are not posterior evidence.

A multi-worker token should contain at least

\[
(\texttt{tree\_id},
\texttt{tree\_generation},
\texttt{request\_id},
\texttt{parent\_node},
\texttt{edge\_idx},
\texttt{edge\_version},
\Path).
\]

\(\SubmitTransitions\) may accept rows out of order. A row is installed only if

\[
\texttt{tree\_generation}\text{ is current}
\]

and

\[
\texttt{edge state}=\texttt{Pending(token)}.
\]

Otherwise the row is stale and is discarded.

Completed backups may use local locks on the path being backed up, or a small
critical section per parent node while its incoming child edge and cache are
updated. There is no need for a global tree lock. The correctness requirement is
that every completed multi-worker backup is equivalent to some serial execution
of Algorithm 6.

For final export, \(\FinishSearch\) and \(\MakeTargets\) still require no pending
requests for the requested trees. An implementation may either drain pending
work or cancel it by incrementing the tree generation.

The mathematical equivalence condition is:

\[
\text{only completed submissions may change posterior or categorical values.}
\]

Pending work may affect scheduling, but it may not affect posterior targets,
policy targets, Q targets, value targets, or categorical solve state.

\section*{14. Acceptance checks}

The implementation is consistent with this spec when:

\begin{enumerate}[leftmargin=2em]
  \item Rust owns only the posterior forest and never steps the environment.
  \item Python submits fused transition plus NN-evaluation results.
  \item Node policy, Q, child, and edge arrays are compact and have length equal
  to the number of valid actions.
  \item Terminal rows submit native WDL categorical outcomes.
  \item Categorical WDL values are stored as native tags, not Dirichlet target
  alphas.
  \item Terminal child intrinsic distance is zero, but parent edge distance to a
  terminal child is one.
  \item A known win solves a node immediately and chooses shortest win.
  \item A known loss requires all moves to be losses and chooses longest loss.
  \item A known draw requires all moves to be draws or losses and chooses by
  \(\texttt{draw\_mode}\).
  \item Search stops below categorical solved nodes.
  \item Unsolved posterior caches are full replacement recomputations, not
  additive deltas.
  \item The categorical cache embedding is used only inside unsolved Dirichlet
  cache recomputation and is never exported as a categorical target alpha.
  \item The policy target for unsolved nodes is the posterior probability of
  optimality, not a visit count.
  \item Export preserves target kind tags so JAX can choose Dirichlet KL losses
  or Dirichlet density NLL losses for categorical rows and edges.
  \item Root advancement only promotes existing children and invalidates stale
  pending transition tokens.
  \item The MVP has at most one pending transition request per tree.
  \item The future multi-worker version has at most one pending transition
  request per edge, and pending requests are not posterior evidence.
  \item No old dirty/repair, node-inflight expansion, child-cache absorbed
  version, or global tree-lock parallelism is required for the MVP.
\end{enumerate}

\end{document}
